% year: תשע"ז
% type: מבחן
% semester: ב
% moed: א
% source: exams/מבחנים/תשעז/מועד א תשעז סמסטר ב.pdf
% number: 6
% points: 10
% categories: functions, derivatives

\begin{question}
(10 נק') תהי $f$ פונקציה חד חד ערכית וגזירה. נתון: $f(1)=2$, $f(2)=4$, $f(3)=5$, $f(4)=6$, $f'(1)=0$, $f'(2)=2$, $f'(3)=0$, $f'(4)=3$.
חשבו: $\left(f^{-1}\right)'(4)$
\end{question}

\begin{hint}
לפי משפט הנגזרת של הפונקציה ההפוכה, $(f^{-1})'(y)=\frac{1}{f'(f^{-1}(y))}$. מהו $f^{-1}(4)$?
\end{hint}

\begin{solution}
מכיוון ש-$f(2)=4$ ו-$f$ חד-חד-ערכית, $f^{-1}(4)=2$.
לפי משפט הנגזרת של הפונקציה ההפוכה: אם $f$ גזירה ב-$x_0$, $f'(x_0)\ne0$ (ו-$f$ רציפה וחד-חד-ערכית בסביבה), אז $f^{-1}$ גזירה ב-$y_0=f(x_0)$ ו-
\[ (f^{-1})'(y_0)=\frac{1}{f'(x_0)}. \]
כאן $x_0=2$, $y_0=4$, ו-$f'(2)=2\ne0$, ולכן
\[ \boxed{(f^{-1})'(4)=\frac{1}{f'(2)}=\frac12}. \]
(שאר הנתונים -- $f'(4)=3$ למשל -- הם מסיחים: יש להעריך את $f'$ בנקודה $f^{-1}(4)=2$ ולא בנקודה $4$.)
\end{solution}
